\chapter{Fazit und Ausblick}
\section{Beantwortung der Forschungsfrage}
Die Forschungsfrage dieser Arbeit lautete: \textit{Wie lässt sich eine 
plattformübergreifende Terminverwaltungsanwendung für Tierhalter im 
Kontext von Veterinärmedizinpraxen konzipieren, umsetzen und hinsichtlich 
ihrer Funktionalität und Usability evaluieren?}

Hinsichtlich der Konzeption wurden zunächst die Anforderungen beider 
Zielgruppen analysiert und in funktionale sowie nicht-funktionale 
Anforderungen überführt. Auf dieser Grundlage wurden die 
Softwarearchitektur, das Komponentendiagramm sowie das 
Navigationskonzept entworfen. Mockups veranschaulichten die geplante 
Benutzeroberfläche und ein Aktivitätsdiagramm beschrieb den 
Terminbuchungsprozess. Da das Backend von \textit{petappoint} im Rahmen eines 
vorangegangenen Projektmoduls bereits entwickelt worden war, 
konzentrierte sich die Konzeption ausschließlich auf die 
Präsentationsschicht.

Hinsichtlich der Umsetzung wurde das Frontend als plattformübergreifende 
mobile Anwendung mit React Native und Expo entwickelt. Als 
UI-Komponentenbibliothek kam Gluestack UI zum Einsatz, während 
TanStack Query das Server-State-Management und Zustand den lokalen 
Client-State übernahmen. Die Mehrsprachigkeit wurde mit i18next 
realisiert. Im Rahmen dieser Arbeit wurde ausschließlich die 
Tierhalter-Rolle implementiert, die alle fünf funktionalen 
Anforderungen dieser Rolle vollständig abdeckt.

Hinsichtlich der Evaluation zeigte die durchgeführte Nutzerstudie 
mit 20 Teilnehmenden, dass die entwickelte Anwendung hinsichtlich 
ihrer Funktionalität und Usability überzeugt. Alle Teilnehmenden 
haben sämtliche Aufgaben erfolgreich abgeschlossen. Der 
durchschnittliche SUS-Score von 94,34 entspricht der Note A und 
belegt eine sehr hohe wahrgenommene Benutzerfreundlichkeit. Der 
Plattformvergleich bestätigte, dass die Anwendung auf iOS und 
Android vergleichbar gut funktioniert. Damit kann die 
Forschungsfrage positiv beantwortet werden.

\section{Fazit}
Die vorliegende Arbeit zeigt, dass sich eine fokussierte 
Terminverwaltungsanwendung für Tierhalter im Kontext von 
Veterinärmedizinpraxen sowohl konzipieren als auch als 
plattformübergreifende mobile Anwendung umsetzen lässt. 
Die Evaluation belegt, dass die entwickelte Lösung die 
definierten Anforderungen der Tierhalter-Rolle erfüllt und 
von den Nutzenden als intuitiv und gebrauchstauglich bewertet 
wird. Der durchschnittliche SUS-Score von 94,34 liegt deutlich oberhalb der Note-A-Grenze und belegt eine hohe wahrgenommene Benutzerfreundlichkeit.

Eine der zentralen Herausforderungen dieser Arbeit bestand darin, 
den Implementierungsumfang bewusst einzugrenzen und sich auf die 
Tierhalter-Rolle zu beschränken, anstatt den Funktionsumfang 
frühzeitig zu erweitern. Diese Eingrenzung des Implementierungsumfangs ermöglichte eine solide und getestete Grundlage für die Weiterentwicklung der Anwendung.

Damit leistet diese Arbeit einen Beitrag zur Digitalisierung 
der Terminverwaltung im veterinärmedizinischen Bereich und 
zeigt, dass der Einsatz moderner plattformübergreifender 
Technologien wie React Native eine effiziente Umsetzung 
solcher Lösungen ermöglicht. Ein wesentlicher nächster Schritt wäre die Erweiterung um die Tierarztpraxis-Rolle, um \textit{petappoint} als vollständige Terminverwaltungslösung einsetzen zu können.

\section{Ausblick}
Die vorliegende Arbeit bildet eine solide Grundlage für die
Weiterentwicklung von \textit{petappoint}. Zunächst sollten die im Rahmen 
der Evaluation gemeldeten Fehler behoben werden, insbesondere die 
fehlende Validierung des Geburtsdatums, die eingeschränkte 
Bearbeitbarkeit von der Startseite sowie das Auftreten 
unerwarteter Termine.

Neben der Frontend-Weiterentwicklung sollte auch die bestehende 
Backend-Infrastruktur überarbeitet werden. Wie in Kapitel~3 
dargelegt, weist das Backend mehrere Schwachstellen auf, die 
bei einer Weiterentwicklung berücksichtigt werden sollten. 
Dazu zählen die unklare Rollentrennung zwischen Tierärzten 
und Personen, die fehlende Rollenhierarchie innerhalb einer 
Tierarztpraxis sowie das Fehlen fachspezifischer Attribute 
für Tierärzte. Darüber hinaus fehlt ein Statusattribut für 
Termine, das eine eindeutige Unterscheidung zwischen freien, 
gebuchten und stornierten Terminen ermöglicht, sowie eine 
dedizierte Tabelle für die Öffnungszeiten der Praxen, die 
für eine automatische Validierung von Buchungen innerhalb 
der Öffnungszeiten erforderlich wäre.

Ein wesentlicher nächster Schritt ist die Implementierung der 
Tierarztpraxis-Rolle, um die Anforderungen FA\#06 und FA\#07 zu 
erfüllen und damit eine vollständige Nutzung der Anwendung durch 
beide Rollen zu ermöglichen. Erst mit der Umsetzung beider Rollen 
kann \textit{petappoint} als vollständige Terminverwaltungslösung eingesetzt werden.

Darüber hinaus bietet die Terminverwaltung weiteres Ausbaupotenzial. 
So sollten Stornierungsfristen eingeführt werden, ab denen eine 
Absage nicht mehr über die App möglich ist. Ebenso wäre ein 
Verschiebungslimit sinnvoll, das verhindert, dass Termine beliebig 
oft umgebucht werden. Ergänzend könnte ein Strike-System eingeführt 
werden, bei dem Tierhalter nach wiederholten kurzfristigen 
Stornierungen vorübergehend keine neuen Termine bei der betroffenen 
Praxis buchen können. Diese Maßnahmen würden die Verbindlichkeit 
der Buchungen erhöhen und den Praxisalltag entlasten.

Hinsichtlich funktionaler Erweiterungen wünschen sich Nutzende 
unter anderem eine Notizfunktion beim Termin, um den Grund des 
Besuchs im Vorfeld zu kommunizieren, die Möglichkeit 
Tierversicherungen zu hinterlegen, erweiterte Tierspezifikationen 
sowie ein Bewertungssystem für Tierarztpraxen. Zusätzlich sollten 
die Suchfilter um Öffnungszeiten am Wochenende erweitert werden, 
da dies von mehreren Teilnehmenden der Nutzerstudie als 
Verbesserungswunsch genannt wurde.

Aus technischer Sicht stellt die Integration von 
Push-Benachrichtigungen direkt auf das Smartphone eine sinnvolle 
Erweiterung dar, um Tierhalter automatisch an bevorstehende Termine 
zu erinnern und über Buchungsbestätigungen oder Stornierungen 
zu informieren. Dies würde die bereits im Backend vorhandene 
E-Mail-Benachrichtigung durch einen direkteren Kommunikationskanal 
ergänzen.

Abschließend bietet auch die gestalterische Weiterentwicklung
der Anwendung Potenzial, etwa durch die Möglichkeit individuelle
Farben für Tierprofile auszuwählen sowie eine überarbeitete visuelle
Gestaltung der Benutzeroberfläche, um die Zufriedenheit der Nutzenden
weiter zu steigern.